\subsection{The main theorems}
\label{series:V:main-theorems}

The resolved fibre product uses the fixed standard
$(3,4,\infty)$ threefold data and the transverse elliptic factor in
\cref{maps:src:E.1.2}. The rank-six construction uses the displayed
integral monodromies of orders five and four, the primitive maximal
cusp and its marked Kuhn fan in \cref{maps:src:F.2.5}. We retain the
complex orientation and column homology vectors. The family parameter
$\xi$ uses the algebraic splitting $\sigma_b$ of
\cref{night:P1.18}; a change of splitting is transported by the full
complex constant stated there. The shifted products in the first part
also include the complex fivefolds explicitly constructed in
\cref{maps:src:E.2.6,maps:src:E.2.14}.

\hypertarget{main-V-A}{}
\paragraph{Theorem A (The resolved transverse elliptic fibre product).}
\label{series:V:A}
The explicit crepant small resolution of the transverse elliptic
fibre product gives a compact complex fourfold with
$K=\Pi^*L^{-5}$ and algebraic dimension two; its Todd and
Pontryagin calculations give signature two and middle form
$\langle1\rangle\oplus\langle1\rangle$.
The computed values $b_1=0$ and $h^{0,1}=1$, preserved by a
K\"ahler modification, obstruct membership in Fujiki's class
$\mathcal C$ through the K\"ahler equality $b_1=2h^{0,1}$
(\SeriesRef{maps:src:E.1.2}, \SeriesRef{maps:src:E.2.1}).

\textit{Mechanism.} The explicit period comparison and a small crepant Kuhn resolution preserve the root canonical calculation and the integral characteristic pairings.
\textit{Relation to literature.}
The threefold factor is the construction of
\cite[Theorems~7.17 and 7.22, pp.~57 and 60]{AC26}, and the
elliptic equation is the one in \cite[\S2, initial construction]{E26}.
The local Kuhn quotient is the $n=3$, $D=C_3$ specialization of
\cite[Theorem~8.7 and Corollary~8.8, attached-source edition, pp.~57--59]{Yang2026}.
The additional global construction is the small crepant resolution
with the integral period morphism
$AZ_3B^{-1}=\operatorname{diag}(Z_2,s+c_E)$ in
\cref{maps:src:E.1.2}; its invariants are calculated in
\cref{maps:src:E.2.1}.  The characteristic polynomials are the
known identities \cite[equations~(1.2)--(1.4), pp.~3--4]{ffEsoleKang},
and the class-$\mathcal C$ obstruction uses
\cite[Theorem~1.1 and Remark~1.2, p.~2]{ffJiaMeng}.
\hyperlink{outline-V-A}{Proof in outline}; full statements: \SeriesRef{maps:src:E.1.2}, \SeriesRef{maps:src:E.2.1}.

\hypertarget{main-V-B}{}
\paragraph{Theorem B (The integral rank-six construction).}
\label{series:V:B}
The displayed rank-six monodromy and the specified affine
fillings produce a three-connected eight-manifold with integral middle
form $3\langle1\rangle\oplus\langle-1\rangle$.  The invariant normal closure of the cusp angular lattice
kills five fibre generators, while the actual order-four boundary
map kills each remaining order-two torsion class with coefficient one.

\textit{Mechanism.} The actual order-four boundary map kills the remaining torsion with coefficient one, and the integral Leray filtration gives the middle lattice.
\textit{Relation to literature.}
The local toroidal filling and its Euler number are the marked
$n=3$ specialization of \cite[Theorem~8.7 and Corollary~8.8,
attached-source edition, pp.~57--59]{Yang2026}.  The global integral specialization
method extends the calculation in \cite[\S7.7]{AC26}.
Here the displayed rank-six restriction maps and order-four
norm/prism map supply the integral groups and the form $I_{3,1}$.
These specific lattice and boundary maps are the additional
calculation asserted in this theorem.
\hyperlink{outline-V-B}{Proof in outline}; full statements: \SeriesRef{maps:src:F.2.5}, \SeriesRef{maps:src:F.2.6}, \SeriesRef{night:P1.13}.

\hypertarget{main-V-C}{}
\paragraph{Theorem C (The genus-two period construction).}
\label{series:V:C}
The pencil $y^2=x^5+\lambda^3x+\lambda^4$ has one nodal
fibre away from $0,\infty$, at $\lambda_*=-256/3125$.
Its good root charts have orders eight and five, and on that root
curve its Hodge bundle is
\[
 \mathcal H=T^3\oplus T,\qquad
 T=\OO((2/5)[\infty]-(3/8)[0]),\qquad \deg T=1/40.
\]
The meromorphic differential $dx/y$ generates $\ell_0=T^3$.
In the full rank-four de Rham bundle, projection of
$\ell_0\wedge\mathcal H_{\rm dR}$ to
$\bigwedge^2\mathcal H_{\rm dR}/\mathbb C\omega$
is a rank-three bundle $P$.  Its holomorphic lift to the specified
rank-six extension is $(\sigma_b+\xi\alpha e)(P)$, with the
normalization of $e,\alpha,\sigma_b$ in \cref{night:P1.15}.
It is transverse to its conjugate on the completed base exactly
for the parameters $\xi\in\mathcal U$ of Theorem~E.
The actual exterior-square and lift maps descend the finite root
orders to five and four; the cusp has period
$Z(q)=(\log q/(2\pi i))I_3+C(q)$ with $C$ holomorphic at zero
(\cref{night:P1.4,night:P1.5,maps:src:F.3.5,maps:src:F.3.6,night:P1.18}).

\textit{Mechanism.} The wedge image of the genus-two Hodge line has rank three, and its split lift becomes transverse after the normalized imaginary shift.
\textit{Relation to literature.}
The construction uses the explicit morphism
$\ell_0\otimes(\mathcal H_{\mathrm{dR}}/\ell_0)
\to\bigwedge^2\mathcal H_{\mathrm{dR}}/\mathbb C\omega$,
$u\otimes[v]\mapsto[u\wedge v]$, and the index-two integral
map in \cref{night:P1.2}.  These formulas specify the
relation to the polarized genus-two variation.  The cited
toroidal quotient is applied through \cref{night:P1.9};
the particular pencil, integral quotient, and split rank-three
plane are the constructions proved here.
\hyperlink{outline-V-C}{Proof in outline}; full statements: \SeriesRef{night:P1.4}, \SeriesRef{night:P1.5}, \SeriesRef{maps:src:F.3.5}, \SeriesRef{maps:src:F.3.6}.

\hypertarget{main-V-D}{}
\paragraph{Theorem D (The oriented diffeomorphism type).}
\label{series:V:D}
The period planes extend over the two finite root charts and
the marked cusp, giving a compact complex fourfold diffeomorphic to
$3\mathbb{HP}^2\#(-\mathbb{HP}^2)$.  Its characteristic vector is
carried explicitly to $(1,1,1,1)$, after which the known smooth
classification and homotopy-sphere absorption identify the oriented
manifold.

\textit{Mechanism.} The middle-lattice isometry transports both Euler and Pontryagin clutching data, after which the cited absorption theorem removes the homotopy sphere.
\textit{Relation to literature.}
The oriented extended-form classification and Kramer--Stolz
absorption theorem are the known results cited in
\cref{maps:src:F.4.1}.  The almost-complex compatibility is
also known from \cite[Theorems~1--2, pp.~2--3]{Yang2012AC}
and \cite[Theorem~1.3 and Corollary~1.4, pp.~3--4]{ffMandal}.
Explicitly, transporting the constructed integrable structure
by the oriented diffeomorphism gives a complex vector bundle
$(TN,J)$ whose realification is $TN$, with
$c_2=-\sum x_i$, $c_4=6u$, and $p_2=14u$.
In Mandal's normalization $x_i^2=\epsilon_i u$,
$\epsilon=(1,1,1,-1)$, his two conditions are
$14\equiv2\pmod4$ and $14-2=2\cdot6$.
The theorem here supplies the compact integrable realization
and its explicit normal-data comparison.
\hyperlink{outline-V-D}{Proof in outline}; full statements: \SeriesRef{maps:src:F.3.7}, \SeriesRef{maps:src:F.4.1}.

\hypertarget{main-V-E}{}
\paragraph{Theorem E (The exact family and its two ends).}
\label{series:V:E}
Put
\[
 c_*=-h_{b,\min}=\frac{96}{625\pi}\operatorname{Re}(XH_\infty)_{03},
 \qquad \mathcal U=\{\xi\in\C:\Im\xi>c_*\},
\]
where $X$ is the normalized integral intertwiner of
\cref{night:P1.15} and $H_\infty$ is its regular infinity-root
fundamental-solution limit, normalized at $\lambda=1$ as in
\cref{night:P1.18}. Thus
$c_*=1.0945402503114027801162682630549457037\ldots$.
The point-normalized parameter is transported by the full complex
translation $\xi_b=\xi_{\rm original}-a_{\rm split}$ of that statement.
The marked family is parametrized exactly by $\mathcal U$:
its members are pairwise nonisomorphic, have algebraic dimension one
and automorphism group $\mathbb C\times\mathbb Z/2$, and admit no
antiholomorphic isomorphism for this fan.  The real-intersection
function determines the parameter domain; the invariant line and
one half-period determine the automorphisms, while the actual
reflected cusp map fails to preserve a ray of the chosen fan.
On the effective root curve put
$L=\OO((4/5)[\infty]-(3/4)[0])$. Then
\[
 R^qh_*\OO=\bigwedge^q(\OO\oplus L^{-1}\oplus L^{-2}),
 \qquad K_M=h^*L^{-6}.
\]
With $p$ indexing rows and $q$ indexing columns, the Hodge table is
\[
\begin{array}{c|ccccc}
h^{p,q}&0&1&2&3&4\\\hline
0&1&1&0&0&0\\
1&0&2&1&0&0\\
2&0&1&6&1&0\\
3&0&0&1&2&0\\
4&0&0&0&1&1
\end{array}
\]
The tangent tuple is $h^q(T_M)=(1,1,1,2,0)$ and its Kuranishi germ
is smooth of dimension one. A full parameter disc across the
transversality boundary has the proper-filling obstruction of
\cref{night:P1A.5}. A proper smooth family on a punctured disc
whose fibres belong to $\mathcal U$ has the smooth filling after
the explicit degree-two torsor cover of \cref{night:P1A.6}.


\textit{Mechanism.} The normalized Frobenius period minimum fixes the exact domain, and the marked period and cap maps determine the family invariants and isomorphisms.
\textit{Relation to literature.}
The root-base and direct-image method is the one in
\cite[Lemmas~6.1--6.2 and Proposition~6.10]{DGKX2026v1};
\cref{night:P1A.1} applies its cup-product and base-change
maps to the displayed rank-three bundle on the $(5,4)$ root curve.
The coefficient $c_*$, the rank-six cohomology and
isomorphism calculations, and the two endpoint assertions are
the specific results proved here.  The torus-deformation input
for \cref{night:P1A.5} is \cite[Theorem~2.1, p.~264]{Catanese2004};
its application is to the smooth compact torus fibres produced
inside that argument.  The punctured-disc filling of
\cref{night:P1A.6} uses its interior parameter limit and
the explicit $\mathbb C\times\mathbb Z/2$ torsor reduction.
The full splitting translation
$\xi_b=\xi_{\rm original}-a_{\rm split}$ is retained.
\hyperlink{outline-V-E}{Proof in outline}; full statements: \SeriesRef{night:P1.18}, \SeriesRef{night:P1A.1}, \SeriesRef{night:P1A.2}, \SeriesRef{night:P1A.3}, \SeriesRef{night:P1A.4}, \SeriesRef{night:P1A.5}, \SeriesRef{night:P1.17}, \SeriesRef{night:P1A.6}.

\hypertarget{main-V-F}{}
\paragraph{Theorem F (The integral input classification).}
\label{series:V:F}
Let $k\in\{2,3\}$ and $\Lambda=\mathbb Z^6$.
Consider ordered integral tuples $(M_1,\ldots,M_k,M_0)$ with each
$M_i$ of finite order and having a nonzero fixed vector, with
$M_0=(M_1\cdots M_k)^{-1}$, and with
$N=M_0-I$ satisfying $N^2=0$, rank three, and
$\operatorname{im}N=\ker N$ as integral sublattices of $\Lambda$.
Every such tuple has a computable canonical representative and an
explicit simultaneous unimodular conjugator; equality of the
representatives is equivalent to simultaneous integral conjugacy.
For input entry height $H$, the finite stopping bound is
$(2H+1)^{36(k+1)}$ candidates before the proved rejection filters.
The ordered sector with both finite linear orders two has exactly
56 classes.  Each has free affine lifts of actual order four,
and every free affine filling of this sector has nontrivial
fundamental group.

Specify the actual free affine actions
$\tau_i(x)=M_ix+v_i/m_i$, with $v_i\in\Lambda$ and actual
order $m_i$, an integral translation clutch $\kappa\in\Lambda$,
and a primitive Kuhn cusp marking.  The algorithms compute the
full fundamental group and every integral homology group of the
resulting marked gluing for every such input.  Their constructive
parametrization and finite tests identify exactly the inputs with
$\pi_1=1$ and $H_2=H_3=0$; there is no fixed bound on the
orders, twists or lattice indices (\cref{maps:src:F.2.39}).

On $B^*=\mathbb P^1\setminus\{p_0,\ldots,p_k\}$ with $k\ge1$,
every representation $\rho:\pi_1(B^*)\to GL^+(6,\mathbb R)$
whose $p_0$ monodromy satisfies $(M_0-I)^2=0$ and
$\operatorname{rank}(M_0-I)=3$ has a holomorphic period family
transverse to its conjugate everywhere on $B^*$. In the marking
$M_0=\left(\begin{smallmatrix}I&I\\0&I\end{smallmatrix}\right)$,
its cusp period is
\[
 Z(t)=\frac{\log t}{2\pi i}I_3+C(t),\qquad C\text{ holomorphic at }0,
\]
with any prescribed finite jet of $C$ at zero. An integral primitive
cusp admits an integral marking. The other $p_i$ remain punctures
in this assertion (\cref{maps:src:F.3.11}).

For a polarized algebraic abelian-surface family with toric-rank-two
cusp, the construction that takes a fibrewise degree-zero line bundle
and quotients its punctured bundle by a lifted translation cannot
have $b_1=b_2=0$ with the stipulated rank-three square-zero Kuhn
cusp; this holds for both proportional and nonproportional source
sections. The filtration and Euler-parity proof is
\cref{maps:src:F.3.4}.

\textit{Mechanism.} Finite-order module restrictions and integral cusp transvections reduce the marked inputs to the stated finite list and its explicit analytic realizations.
\textit{Relation to literature.} The integral and Oka inputs are credited at their exact locators in \cref{maps:src:F.2.39,maps:src:F.3.11}; the statement applies them to every retained marked input.
\hyperlink{outline-V-F}{Proof in outline}; full statements: \SeriesRef{maps:src:F.2.39}, \SeriesRef{maps:src:F.3.11}, \SeriesRef{maps:src:F.3.4}.

\subsection{The proofs in outline}
\hypertarget{outline-V-A}{}
\subsubsection*{Theorem A: The resolved transverse elliptic fibre product}
\begin{enumerate}[label=\arabic*.,leftmargin=*]
\item\hypertarget{step-V-A-1}{} Pull back the Weierstrass family to the root curve of the
threefold fibration.  The integral matrices $A,B$ give the actual
period comparison $AZ_3B^{-1}=\operatorname{diag}(Z_2,s+c_E)$.
\textit{Expansion:} \SeriesRef{maps:src:E.1.2}.
\item\hypertarget{step-V-A-2}{} The unimodular Kuhn fan resolves the local fibre product
without adding rays.  Its period action is proper and its marked
section extends face by face.  The resolution is therefore small
and crepant.
\textit{Expansion:} \SeriesRef{maps:src:E.1.2}.
\item\hypertarget{step-V-A-3}{} The canonical formula on the root curve gives
$K=\Pi^*L^{-5}$.  The split relative elliptic complex and Leray
give $(h^{0,q})=(1,1,0,0,0)$; the explicit meromorphic product
comparison gives algebraic dimension two.
\textit{Expansion:} \SeriesRef{maps:src:E.2.1}.
\item\hypertarget{step-V-A-4}{} A K\"ahler modification in class $\mathcal C$ preserves
$b_1$ and $h^{0,1}$.  Here $b_1=0$ and $h^{0,1}=1$, contradicting
the K\"ahler equality $b_1=2h^{0,1}$.  This proves the stated
exclusion through these invariants.
\textit{Expansion:} \SeriesRef{maps:src:E.2.1}.
\item\hypertarget{step-V-A-5}{} The universal Todd and Pontryagin formulas give
$c_2^2=2$, $c_4=6$ and signature two.  The middle lattice is
positive definite unimodular of rank two, hence has the displayed
form.  The body records the precise scope of each cited formula.
\textit{Expansion:} \SeriesRef{maps:src:E.2.1}.
\end{enumerate}

\hypertarget{outline-V-B}{}
\subsubsection*{Theorem B: The integral rank-six construction}
\begin{enumerate}[label=\arabic*.,leftmargin=*]
\item\hypertarget{step-V-B-1}{} The matrices $M_1=A^t$, $M_2=B^t$ have orders five and four, and
  $N=(M_1M_2)^{-1}-I$ has square zero and primitive image of rank three.
  The displayed unimodular marking takes $N$ to
  $\left(\begin{smallmatrix}0&I_3\\0&0\end{smallmatrix}\right)$;
  the invariant covector $\ell(x)=x_0$ proves freeness of every
  nonidentity affine power
  (\cref{maps:src:F.2.5}, \cref{night:P1.1}).
\textit{Expansion:} \SeriesRef{maps:src:F.2.5}, \SeriesRef{night:P1.1}.
\item\hypertarget{step-V-B-2}{} The cusp and finite quotient charts glue to a closed smooth
  manifold.  Killing the invariant normal closure of the cusp angular lattice
  leaves the central fibre generator $h$ and the relations $5a=h$, $4b=-h$, $a+b=0$;
  their integral relation matrix is unimodular, so the fundamental
  group vanishes (\cref{maps:src:F.2.5}, \cref{night:P1.10}).
\textit{Expansion:} \SeriesRef{maps:src:F.2.5}, \SeriesRef{night:P1.10}.
\item\hypertarget{step-V-B-3}{} Unimodular coordinates identify the finite core monodromies
  with a five-cycle and a four-cycle plus a sign.  Their actual
  covering maps induce invariant vertical cochains and norm sums
  on the prism cochains.  Together with the primitive cusp collapse,
  these maps give the displayed integral sheaf complex; its signed
  orbit bases and determinant minors give the sheaf groups
  (\cref{maps:src:F.2.6}, items~1--4; \cref{night:P1.11}).
\textit{Expansion:} \SeriesRef{maps:src:F.2.6}, \SeriesRef{night:P1.11}.
\item\hypertarget{step-V-B-4}{} For the order-four core and normal loops $g,h$, the fibre and
  transverse loops are $v=g^4h^{-1}$ and $t=g$; their integral
  basis matrix has determinant one.  A torsion representative $w$
  pulls back to $(Nw,w)=(0,w)$ in these two directions.
  Thus its transgression is $\pm[w]$.  The printed modulo-two
  rows detect all three nonzero global classes
  (\cref{maps:src:F.2.6}, item~5;
  \cref{mapsaudit:order-four-torsion-transgression}).
\textit{Expansion:} \SeriesRef{mapsaudit:order-four-torsion-transgression}.
\item\hypertarget{step-V-B-5}{} The three torsion differentials are isomorphisms onto
  $\mathbb Z/2$.  On the cusp term $\mathbb Z\oplus\mathbb Z/2$,
  the remaining kernel is generated by $(1,\epsilon)$, where
  $\epsilon$ is the image of its free generator.
  The Leray filtration therefore gives $H^3=0$ and
  $H^4=\mathbb Z^4$.  Universal coefficients, duality and
  Hurewicz give the asserted free homology and three-connectivity
  (\cref{maps:src:F.2.6}, item~6; \cref{night:P1.12}).
\textit{Expansion:} \SeriesRef{maps:src:F.2.6}, \SeriesRef{night:P1.12}.
\item\hypertarget{step-V-B-6}{} The compatible almost complex structure has $c_1=c_3=0$.
  Todd integrality and the signature bound $|\sigma|\leq4$ force
  $c_2^2=2$, $p_1^2=8$, $p_2=14$ and $\sigma=2$.
  Its unimodular middle lattice is consequently the odd form
  $I_{3,1}$ (\cref{night:P1.13}).
\textit{Expansion:} \SeriesRef{night:P1.13}.
\end{enumerate}

\hypertarget{outline-V-C}{}
\subsubsection*{Theorem C: The genus-two period construction}
\begin{enumerate}[label=\arabic*.,leftmargin=*]
\item\hypertarget{step-V-C-1}{} The pencil has one node at $\lambda_*=-256/3125$ and good
  reductions after root orders eight and five.  Its Hodge bundle on
  that root curve is $T^3\oplus T$, with
  $T=\mathcal O((2/5)[\infty]-(3/8)[0])$ and
  $\deg T=1/40$ (\cref{night:P1.4}).
\textit{Expansion:} \SeriesRef{night:P1.4}.
\item\hypertarget{step-V-C-2}{} In the exterior-square lattice, projection induces
  $j:\omega^\perp\to(\bigwedge^2\mathbb Z^4)/\mathbb Z\omega$
  with Smith factors $(1,1,1,1,2)$.  This exact index-two map identifies
  the rational systems, and the quotient is the dual of the primitive
  integral lattice.  The lower five-dimensional blocks of the given
  rank-six tuple are the quotient blocks under this map
  (\cref{night:P1.2}).
\textit{Expansion:} \SeriesRef{night:P1.2}.
\item\hypertarget{step-V-C-3}{} Put $P=p(\ell_0\wedge\mathcal H_{\mathrm{dR}})$, where
  $\ell_0=\langle dx/y\rangle$ and $p$ is projection modulo $\omega$.
  The wedge map has rank three and its image meets $\mathbb C\omega$
  trivially: a decomposable two-vector has square zero, whereas
  $\omega^2\ne0$.  This proves the required rank before a lift is
  chosen (\cref{night:P1.5}, \cref{maps:src:F.3.5}).
\textit{Expansion:} \SeriesRef{night:P1.5}, \SeriesRef{maps:src:F.3.5}.
\item\hypertarget{step-V-C-4}{} The orbifold extension over $P$ splits because
  $H^1(P^*)=0$, as computed from its two line-bundle extensions and
  their integral floor degrees.  A chosen splitting $\sigma$ and
  the canonical functional $\alpha:P\to\mathcal O$ give
  $F_\xi=(\sigma+\xi\alpha e)P$
  (\cref{maps:src:F.3.5}, \cref{night:P1.7}).
\textit{Expansion:} \SeriesRef{maps:src:F.3.5}, \SeriesRef{night:P1.7}.
\item\hypertarget{step-V-C-5}{} On the real line $P\cap\overline P$, choose $n$ with
  $\alpha(n)=1$ and write
  $\sigma(n)-\overline{\sigma(n)}=2ih e$.
  Then $F_\xi\cap\overline{F_\xi}=0$ exactly when
  $h+\operatorname{Im}\xi\ne0$.  At the cusp the period matrix is
  $Z=(\log q)/(2\pi i)I_3+C(q)$, with $C$ holomorphic and bounded;
  hence $h\to+\infty$ and a sufficiently large imaginary shift
  works on the whole completed base
  (\cref{maps:src:F.3.6}, \cref{night:P1.8}).
\textit{Expansion:} \SeriesRef{maps:src:F.3.6}, \SeriesRef{night:P1.8}.
\end{enumerate}

\hypertarget{outline-V-D}{}
\subsubsection*{Theorem D: The oriented diffeomorphism type}
\begin{enumerate}[label=\arabic*.,leftmargin=*]
\item\hypertarget{step-V-D-1}{} The lifted periods commute with the finite root actions and
  have the prescribed bounded cusp correction, so the same quotient
  and toroidal charts are holomorphic and proper.  Their collar
  identifications are the marked maps used in the smooth construction
  (\cref{maps:src:F.3.7}, \cref{night:P1.9}).
\textit{Expansion:} \SeriesRef{maps:src:F.3.7}, \SeriesRef{night:P1.9}.
\item\hypertarget{step-V-D-2}{} The preceding integral calculation gives the intersection
  lattice $I_{3,1}$.  The spin characteristic vector $p_1/2$ has
  square two, by the Chern-number calculation
  (\cref{night:P1.13}).
\textit{Expansion:} \SeriesRef{night:P1.13}.
\item\hypertarget{step-V-D-3}{} In an odd diagonal basis write its coordinates, after sign
  changes, as positive odd integers $(a,b,c,d)$.  A reflection in
  $(1,1,1,1)$ replaces $d$ by $2d-a-b-c$, whose absolute value is
  smaller when $d>1$.  Iteration terminates at $(1,1,1,1)$, proving
  an isometry of the lattice together with its characteristic vector
  (\cref{maps:src:F.4.1}).
\textit{Expansion:} \SeriesRef{maps:src:F.4.1}.
\item\hypertarget{step-V-D-4}{} In the quaternionic clutching coordinates, the actual map is
  $(a,b)\mapsto(e,p_1/2)=(a-b,a+b)$, with inverse
  $((e+p_1/2)/2,(p_1/2-e)/2)$.  Thus the pairing and characteristic
  vector determine the normal datum; the Poincar\'e-duality map
  transports it with the middle lattice isometry
  (\cref{maps:src:F.4.1}).
\textit{Expansion:} \SeriesRef{maps:src:F.4.1}.
\item\hypertarget{step-V-D-5}{} Each quaternionic projective-plane summand has $p_1=2u$.
  The oriented connected-sum pairing and characteristic vector are
  therefore the same pair $\bigl(I_{3,1},(1,1,1,1)\bigr)$.
  Wall's stated extended-form classification gives the oriented
  diffeomorphism up to a homotopy eight-sphere
  (\cref{maps:src:F.4.1}).
\textit{Expansion:} \SeriesRef{maps:src:F.4.1}.
\item\hypertarget{step-V-D-6}{} Kramer--Stolz's Theorem A absorbs that sphere in one positive
  $\mathbb{HP}^2$ summand.  Its diffeomorphism has degree $+1$ because
  the square of a middle generator evaluates to $1$ on both sides.
  Oriented connected-sum maps therefore give the required oriented
  diffeomorphism of the full fourfold
  (\cref{maps:src:F.4.1}).
\textit{Expansion:} \SeriesRef{maps:src:F.4.1}.
\end{enumerate}

\hypertarget{outline-V-E}{}
\subsubsection*{Theorem E: The exact family and its two ends}
\begin{enumerate}[label=\arabic*.,leftmargin=*]
\item\hypertarget{step-V-E-1}{} The real coordinate and marked deck maps preserve the normalized
  height. The infinity-root rotation fixes its differential at zero;
  an invertible real system and the positive Hessian prove a strict
  local minimum. Convergent root and cusp series and exact collar
  continuation enclose all 1600 leaves of the closed covers above
  that value. The actual cusp graph has a positive determinant on
  the remaining inner disc and tends to positive height without
  bound. These maps prove the unique global minimum and identify
  the transversality domain exactly with $\mathcal U$
  (\cref{night:P1.18}, proof, items 1--7).
\textit{Expansion:} \SeriesRef{night:P1.18}.
\item\hypertarget{step-V-E-2}{} On the effective root curve of orders five and four put
  $L=\mathcal O((4/5)[\infty]-(3/4)[0])$.  The direct images are
  $R^qh_*\mathcal O=\bigwedge^q
  (\mathcal O\oplus L^{-1}\oplus L^{-2})$, and
  $K_M=h^*L^{-6}$.  The toric normalization complex proves the
  exterior-algebra assertion at the cusp, and the root-chart
  invariant cohomology proves it at the finite fibres
  (\cref{night:P1A.1}).
\textit{Expansion:} \SeriesRef{night:P1A.1}.
\item\hypertarget{step-V-E-3}{} The logarithmic cotangent sequence and its explicit residue
  matrices give the Hodge table below.  Serre duality applied to
  the twisted logarithmic and residue sequences gives
  $h^q(T_M)=(1,1,1,2,0)$.  The normal tangent map and the actual
  period lift give a nonzero Kodaira--Spencer class; the local
  analytic-ring comparison makes the Kuranishi germ smooth of
  dimension one (\cref{night:P1A.2}, \cref{night:P1A.3}).
\textit{Expansion:} \SeriesRef{night:P1A.2}, \SeriesRef{night:P1A.3}.
\item\hypertarget{step-V-E-4}{} Restriction of a horizontal divisor to a regular torus fibre
  gives an invariant alternating class.  The computed invariant
  space is zero, while positive pairing with a torus Kähler form
  rules out a nonzero effective divisor of zero class.  Thus the
  meromorphic field is the base field $\mathbb C(\lambda)$.
  Every isomorphism therefore preserves the base and its three
  distinguished values.  The simultaneous linear centralizer is
  $\mathbb Q I$; the order-five affine norm excludes its negative
  integral unit.  The translation sheaf has identity component
  $\mathbb C$ and an integral parabolic quotient $\mathbb Z/2$,
  represented by the displayed half-period, proving the
  automorphism group and equality of parameters for isomorphic
  members (\cref{night:P1A.4}).
\textit{Expansion:} \SeriesRef{night:P1A.4}.
\item\hypertarget{step-V-E-5}{} The two integral anti-intertwiners are $\pm R$.  One is excluded
  by the finite cap norm.  For the other, the actual based cusp
  correction acts on rays by $-P$, and takes $e_2$ to $(0,-1,1)$,
  outside the component's rays $\pm(\{0,1\}^3\setminus\{0\})$.
  The corresponding monomial map contracts a divisor, contradicting
  a biholomorphism (\cref{night:P1.17}).
\textit{Expansion:} \SeriesRef{night:P1.17}.
\item\hypertarget{step-V-E-6}{} A proposed proper Hausdorff filling across a boundary point of
  a full parameter disc recovers the marked base from the moving
  pencil of $K_M^{-4}$, whose two sections have ratio $\lambda$.
  Closed relative forms give actual torus fibres outside a proper
  analytic subset.  Their marked periods continue as $F_\xi$,
  while the real hypersurface $h_b+\operatorname{Im}\xi=0$
  cannot be contained in that analytic subset.  A torus fibre
  there violates real transversality (\cref{night:P1A.5}).  For a family on a punctured disc whose
  fibres belong to $\mathcal U$, the unique classifying map is
  holomorphic; its Cayley transform extends and the maximum principle
  keeps its limit inside $\mathcal U$.  After the explicit degree-two
  removal of the finite torsor, that family has a smooth filling
  (\cref{night:P1A.6}).
\textit{Expansion:} \SeriesRef{night:P1A.5}, \SeriesRef{night:P1A.6}.
\end{enumerate}

\hypertarget{outline-V-F}{}
\subsubsection*{Theorem F: The integral input classification}
\begin{enumerate}[label=\arabic*.,leftmargin=*]
\item\hypertarget{step-V-F-1}{} Primitive square-zero cusp matrices have the explicit integral
  normal form and centralizer
  $\left(\begin{smallmatrix}P&Q\\0&P\end{smallmatrix}\right)$.
  Eigenvalue one is necessary for a free affine lift, leaving the
  finite cyclotomic list and its complex-compatibility filter
  (\cref{maps:src:F.2.1}, \cref{maps:src:F.2.2}).
\textit{Expansion:} \SeriesRef{maps:src:F.2.1}, \SeriesRef{maps:src:F.2.2}.
\item\hypertarget{step-V-F-2}{} The rational intertwiner algebra, its radical and its
  semisimple quotient reduce the simultaneous integral-conjugacy
  problem to the cited lattice-isomorphism algorithm in its verified
  hypothesis.  Searching by height and lexicographic order terminates
  by the input height, yielding a canonical representative and an
  explicit unimodular conjugator (\cref{maps:src:F.2.36}, \cref{maps:src:F.2.37}).
\textit{Expansion:} \SeriesRef{maps:src:F.2.36}, \SeriesRef{maps:src:F.2.37}.
\item\hypertarget{step-V-F-3}{} For that conjugator $R$, the coordinate map $x\mapsto Rx$
  transports $(M_i,v_i,\kappa,P)$ to
  $(RM_iR^{-1},Rv_i,R\kappa,RP)$.  Since
  $(RP)^{-1}(Rx)=P^{-1}x$, it intertwines the specified cusp charts
  as well as the affine caps and clutch.  The resulting finite
  cellular cone and group presentation compute all integral homology
  and the full fundamental group (\cref{maps:src:F.2.31}, \cref{maps:src:F.2.35}, \cref{maps:src:F.2.39}).
\textit{Expansion:} \SeriesRef{maps:src:F.2.31}, \SeriesRef{maps:src:F.2.35}, \SeriesRef{maps:src:F.2.39}.
\item\hypertarget{step-V-F-4}{} The exact two-point survivor tests and the three-point
  orbifold quotient give the complete constructive parametrization;
  the separate involution branch has 56 classes.  The displayed
  unbounded trace family proves why the terminating bound depends on
  the supplied input (\cref{maps:src:F.2.30}, \cref{maps:src:F.2.38}, \cref{maps:src:F.2.39}).
\textit{Expansion:} \SeriesRef{maps:src:F.2.30}, \SeriesRef{maps:src:F.2.38}, \SeriesRef{maps:src:F.2.39}.
\item\hypertarget{step-V-F-5}{} Nine explicitly given complex-line directions and a uniform
  singular-value bound put a spanning tube in the unpolarized period
  domain.  The verified local-spray approximation theorem, with the
  relative extension-point argument proved here, produces a
  holomorphic transverse period section and retains any prescribed
  bounded cusp jet (\cref{maps:src:F.3.8}, \cref{maps:src:F.3.10}, \cref{maps:src:F.3.11}).
\textit{Expansion:} \SeriesRef{maps:src:F.3.8}, \SeriesRef{maps:src:F.3.10}, \SeriesRef{maps:src:F.3.11}.
\item\hypertarget{step-V-F-6}{} For a polarized abelian-surface source, semisimplicity and the
  maximal square-zero cusp force the middle rank-four system to have
  no invariant or coinvariant vectors.  The actual filtration
  $0\subset\mathbb Q\subset E\subset T$, with middle quotient that
  system and top quotient $\mathbb Q$, then forces both relevant
  invariant ranks to be one; the integral Euler-parity relation
  excludes $b_1=b_2=0$ (\cref{maps:src:F.3.4}).
\textit{Expansion:} \SeriesRef{maps:src:F.3.4}.
\end{enumerate}

